% Lösungen Einheit 3 von 4 – Baumdiagramm und Pfadregeln (zusammenbau v0.1)
\einheitenkopf{Einheit 3 von 4 · Baumdiagramm und Pfadregeln}

\erg{17}{a) 2 Pfade: (rot; blau), (blau; rot)}
\erg{18}{a) Gegenereignis: keine Sechs in beiden Würfen}
\erg{19}{a) ohne Zurücklegen \quad b) an jedem Ast $\frac{1}{2}$ \quad c) grüner Würfel ungerade $\frac{3}{6}$; gelber Würfel an beiden Punkten gerade $\frac{5}{6}$, ungerade $\frac{1}{6}$ \quad d) $\frac{3}{8} \cdot \frac{3}{8} = \frac{9}{64}$ \quad e) $\frac{3}{4} \cdot \frac{2}{4} + \frac{1}{4} \cdot \frac{1}{4} = \frac{6}{16} + \frac{1}{16} = \frac{7}{16}$ \quad f) $\frac{25}{64} + \frac{4}{64} + \frac{1}{64} = \frac{30}{64} = \frac{15}{32}$ \quad g) $1 - \frac{5}{6} \cdot \frac{5}{6} = 1 - \frac{25}{36} = \frac{11}{36}$ \quad h) $1 - \frac{3}{6} \cdot \frac{5}{6} = 1 - \frac{15}{36} = \frac{21}{36} = \frac{7}{12}$ \quad i) $\frac{1}{2} \cdot \frac{1}{2} \cdot \frac{1}{2} = \frac{1}{8}$ \quad j) Nein: 66 hat $\frac{1}{4} \cdot \frac{2}{4} = \frac{2}{16}$, 77 hat $\frac{1}{4} \cdot \frac{1}{4} = \frac{1}{16}$. \quad k) z. B. links 5, 5, 5, 6 und rechts 7, 7, 8, 8: $\frac{3}{4} \cdot \frac{2}{4} = \frac{6}{16} = \frac{3}{8}$ \quad l) an jedem Ast Stern $\frac{1}{4}$, leer $\frac{3}{4}$; $3 \cdot \frac{1}{4} \cdot \frac{1}{4} \cdot \frac{3}{4} + \frac{1}{4} \cdot \frac{1}{4} \cdot \frac{1}{4} = \frac{9}{64} + \frac{1}{64} = \frac{10}{64} = \frac{5}{32}$}

19b)

\baumzwei{Kopf/\frac{1}{2},Zahl/\frac{1}{2}}{Kopf/\frac{1}{2},Zahl/\frac{1}{2}}{Kopf/\frac{1}{2},Zahl/\frac{1}{2}}


19c)

\baumzwei{gerade/\frac{3}{6},ungerade/\frac{3}{6}}{gerade/\frac{5}{6},ungerade/\frac{1}{6}}{gerade/\frac{5}{6},ungerade/\frac{1}{6}}


19k)

\kreisdiagramm{5/1,5/1,5/1,6/1} \kreisdiagramm{7/1,7/1,8/1,8/1}


19l)

\baumdreigleich{Stern/\frac{1}{4},leer/\frac{3}{4}}

\erg{20}{a) Fehler: Die beiden Pfade werden multipliziert statt addiert. Richtig: $\frac{1}{4} \cdot \frac{3}{4} + \frac{3}{4} \cdot \frac{1}{4} = \frac{6}{16} = \frac{3}{8}$. \quad b) Die Äste an einem Punkt zeigen alle möglichen Ergebnisse dieser Stufe; genau eines davon tritt sicher ein. \quad c) z. B. ein Glücksrad mit fünf gleich großen Feldern, zwei rot und drei blau, zweimal gedreht \quad d) $0{,}8 \cdot 0{,}8 = 0{,}64$}
